\section{Fourier transforms and delta functions}\label{app:C.1}

Fermi's golden rule rests on one limit: the function $\sin^2(\xi t/2)/\xi^2$ of \cref{eq:2.1-Pstatic} grows in height as $t^2$ and shrinks in width as $1/t$, and at long times it acts on smooth functions as $(\pi t/2)\,\delta(\xi)$ (\cref{eq:2.1-deltalimit}). \Cref{sec:2.1} used the limit and the integral $\int\dd u\,\sin^2u/u^2 = \pi$ behind it; \cref{sec:2.2} used the rule for a delta function of a function. This section proves the limit, the integral and the rule for a delta function of a function, together with the Fourier identities that the rest of the book uses, in the conventions of \cref{app:A.3}.

\subsection{Delta functions as limits}

A delta function is defined by what it does under an integral. For every test function $w(x)$, a smooth function that falls off fast enough for the integrals below to converge,
\begin{equation}
  \int_{-\infty}^{\infty}\dd x\;\delta(x)\,w(x) = w(0) .
  \label{eq:C.1-delta}
\end{equation}
No ordinary function does this; $\delta$ is a distribution, a linear rule that assigns a number to each test function~\cite[Sect.~1.16]{olver2010nist}. In physics it always appears as the limit of ordinary functions that become tall and narrow, and the useful question is when and how fast the limit is reached.

The basic construction takes any integrable $f(u)$ with $\int\dd u\,f(u) = 1$ and squeezes it, $f_t(x) = t\,f(tx)$. The area of $f_t$ stays one while its width shrinks as $1/t$. With $u = tx$,
\begin{equation*}
  \int\dd x\,f_t(x)\,w(x) = \int\dd u\,f(u)\,w(u/t) \;\longrightarrow\; w(0)\int\dd u\,f(u) = w(0) \qquad (t\to\infty),
\end{equation*}
because $w(u/t) \to w(0)$ at every fixed $u$, and $|f(u)\,w(u/t)| \le |f(u)|\max|w|$ lets the limit pass under the integral (dominated convergence). The family $f_t$ therefore tends to $\delta$. Three members recur in this book: the Gaussian $f(u) = \ee^{-u^2/2}/\sqrt{2\pi}$, which describes a broadened line on a computational grid (\cref{app:C.2}); the Lorentzian $f(u) = \pi^{-1}/(1 + u^2)$, which describes a decaying state (\cref{sec:2.4,app:C.4}); and the golden-rule function, $f(u) = (2/\pi)\sin^2(u/2)/u^2$.

How fast the limit is reached depends on the tails of $f$. For an even $f$ with a finite second moment $\langle u^2\rangle_f$, expanding $w(u/t)$ in powers of $u/t$ gives $w(0) + \tfrac12 w''(0)\langle u^2\rangle_f/t^2$, and the correction to $w(0)$ is of relative order $1/t^2$. The Gaussian is of this kind. The Lorentzian and the golden-rule function fall off only as $1/u^2$, and their second moments diverge. The expansion then fails: their corrections are of order $1/t$, and they are set by the values of $w$ far from the origin, not by its curvature at the origin. \Cref{eq:C.1-sincorr} below gives the correction for the golden-rule function and \cref{app:C.4} the one for the Lorentzian.

\subsection{The golden-rule limit}

The golden-rule limit needs two integrals, the Dirichlet integral $\int_0^\infty\dd u\,\sin u/u = \pi/2$ and $\int\dd u\,\sin^2u/u^2 = \pi$. The Dirichlet integral follows from a convergence factor. Since $\int_0^\infty\dd u\,\ee^{-su}\sin u = 1/(1 + s^2)$ for $s > 0$, and $1/u = \int_0^\infty\dd s\,\ee^{-su}$ for $u > 0$,
\begin{equation*}
  \int_0^\infty\dd u\,\frac{\sin u}{u} = \int_0^\infty\dd s\int_0^\infty\dd u\,\ee^{-su}\sin u = \int_0^\infty\frac{\dd s}{1 + s^2} = \frac{\pi}{2} .
\end{equation*}
The exchange of the two integrals is legitimate because the double integral converges absolutely once a factor $\ee^{-s_0u}$ with $s_0 > 0$ is included, and the result is continuous as $s_0 \to 0$. The integrand is even, so $\int_{-\infty}^\infty\dd u\,\sin u/u = \pi$. Integrating by parts, with $\sin^2u/u \to 0$ at both ends and $\dd(\sin^2u)/\dd u = \sin 2u$,
\begin{equation}
  \int_{-\infty}^{\infty}\dd u\,\frac{\sin^2u}{u^2} = \Bigl[-\frac{\sin^2u}{u}\Bigr]_{-\infty}^{\infty} + \int_{-\infty}^{\infty}\dd u\,\frac{\sin 2u}{u} = \pi ,
  \label{eq:C.1-sinsq}
\end{equation}
where the last integral is the Dirichlet integral with $u \to 2u$. \Cref{eq:C.1-sinsq} is the integral quoted in \cref{sec:2.1}. The central lobe of the integrand, $|u| < \pi$, holds \SI{90.3}{\percent} of the integral.

Write $\sin^2(\xi t/2) = (1 - \cos\xi t)/2$ and split off $w(0)$:
\begin{equation*}
  \int\dd\xi\,\frac{\sin^2(\xi t/2)}{\xi^2}\,w(\xi)
  = w(0)\,\frac{\pi t}{2} + \frac12\int\dd\xi\,\frac{w(\xi) - w(0)}{\xi^2} - \frac12\int\dd\xi\,\cos(\xi t)\,\frac{w(\xi) - w(0)}{\xi^2} .
\end{equation*}
The first term uses \cref{eq:C.1-sinsq} with $u = \xi t/2$. In the other two, $w(\xi) - w(0) = w'(0)\xi + O(\xi^2)$ near the origin. The odd part $w'(0)\xi$ makes the integrand odd and singular as $1/\xi$, and the contributions from the two sides of $\xi = 0$ cancel if the integral is taken symmetrically about $\xi = 0$; the $O(\xi^2)$ term leaves an integrable integrand. The second term is then a constant, the Hadamard finite part of $\int\dd\xi\,w(\xi)/\xi^2$ defined in \cref{eq:2.4-fp}. The third term tends to zero as $t \to \infty$: the integral of a rapidly oscillating cosine against an integrable function vanishes in the limit (the Riemann--Lebesgue lemma~\cite[Eq.~(1.8.10)]{olver2010nist}), and the singular odd part contributes nothing by symmetry. Hence
\begin{equation}
  \int\dd\xi\,\frac{\sin^2(\xi t/2)}{\xi^2}\,w(\xi) = \frac{\pi t}{2}\,w(0) + \frac12\operatorname{fp}\!\int\dd\xi\,\frac{w(\xi)}{\xi^2} + o(1) \qquad (t\to\infty),
  \label{eq:C.1-sincorr}
\end{equation}
where $o(1)$ stands for terms that vanish as $t\to\infty$. Dividing by $\pi t/2$ proves \cref{eq:2.1-deltalimit}. For a $w$ that varies on an energy scale $\delta_w$, the finite part is of size $w(0)/\delta_w$, so the correction relative to the leading term is of order $1/(t\,\delta_w)$, the estimate quoted in \cref{sec:2.1}. A constant $w$ has finite part zero, and then the area $\pi t/2$ is exact at every $t$.

The time derivative of \cref{eq:2.1-Pstatic} brings in $\sin(\xi t)/\xi$, which has area $\pi$ at every $t > 0$ by the Dirichlet integral. The function $\sin(\xi t)/\xi$ is not absolutely integrable, so the dominated-convergence argument does not apply, but splitting off $w(0)$ as above still works. With $w(\xi) = w(0) + [w(\xi) - w(0)]$, the first part gives $\pi w(0)$. The second gives $\int\dd\xi\,\sin(\xi t)\,[w(\xi) - w(0)]/\xi$, whose integrand is a sine times a function that is finite at the origin. This integral vanishes as $t\to\infty$ by the Riemann--Lebesgue lemma (at large $|\xi|$, where $w(0)/\xi$ is not integrable, the oscillation of the sine still makes the contribution vanish). Hence
\begin{equation}
  \frac{\sin\xi t}{\xi} \;\longrightarrow\; \pi\,\delta(\xi) \qquad (t\to\infty) .
  \label{eq:C.1-sinc}
\end{equation}
The convergence is weaker than for $\sin^2(\xi t/2)/\xi^2$. The function $\sin(\xi t)/\xi$ changes sign, its absolute area grows as $\ln t$, and the limit holds only under an integral against a smooth function. Because of this weaker convergence, \cref{sec:2.1} reads the rate off the probability, whose kernel is positive, as well as off the time derivative of the probability.

\subsection{Delta functions of functions}

Let $g(x)$ be smooth with simple roots $x_i$, that is $g(x_i) = 0$ and $g'(x_i) \neq 0$. Near each root $g(x) \approx g'(x_i)(x - x_i)$, and a delta function of $g$ is supported only at the roots. Changing the variable to $y = g(x)$ in a small interval around $x_i$, with $\dd x = \dd y/|g'(x)|$ (the absolute value appears because the orientation of the interval is restored after the substitution),
\begin{equation*}
  \int_{x_i - \epsilon}^{x_i + \epsilon}\dd x\,\delta(g(x))\,w(x) = \int\dd y\,\delta(y)\,\frac{w(x(y))}{|g'(x(y))|} = \frac{w(x_i)}{|g'(x_i)|} .
\end{equation*}
Summing over the roots,
\begin{equation}
  \delta(g(x)) = \sum_i\frac{\delta(x - x_i)}{|g'(x_i)|} .
  \label{eq:C.1-deltag}
\end{equation}
Two cases appear in \cref{chap:2}. The scaling rule $\delta(cx) = \delta(x)/|c|$ is the case of one root. Energy conservation in elastic scattering, $\delta(p'^2/2\mu_{\chi e} - p^2/2\mu_{\chi e})$ with $p, p' \ge 0$, has the root $p' = p$ with $g'(p) = p/\mu_{\chi e}$ and gives $(\mu_{\chi e}/p)\,\delta(p' - p)$, as used in \cref{sec:2.2}. A double root, where $g'(x_i) = 0$, makes the right side of \cref{eq:C.1-deltag} infinite; it signals that the delta function must be kept inside a larger integral, as happens at the van Hove singularities of a density of states (\cref{app:C.2}).

In three dimensions $\delta^3(\vect{x}) = \delta(x_1)\delta(x_2)\delta(x_3)$. The analogue of \cref{eq:C.1-deltag} for a scalar function of a vector is an integral over the surface on which the function vanishes,
\begin{equation}
  \int\dd^3k\,\delta(g(\vect{k}))\,w(\vect{k}) = \int_{g(\vect{k}) = 0}\frac{\dd S}{|\nabla g(\vect{k})|}\,w(\vect{k}) ,
  \label{eq:C.1-surface}
\end{equation}
which follows from \cref{eq:C.1-deltag} applied along the normal to the surface at each point, with $\dd S$ the element of area. \Cref{app:C.2} uses \cref{eq:C.1-surface} for densities of states.

\subsection{Fourier identities}

The transforms of \cref{eq:A.3-space,eq:A.3-time} are inverse to each other because the plane waves are complete,
\begin{equation}
  \int_{-\infty}^{\infty}\dd t\,\ee^{\ii\omega t} = 2\pi\,\delta(\omega), \qquad
  \int\dd^3r\,\ee^{\ii\vect{q}\cdot\vect{r}} = (2\pi)^3\,\delta^3(\vect{q}) .
  \label{eq:C.1-complete}
\end{equation}
The first identity is the golden-rule limit once more. The integral over a finite window, $\int_{-T/2}^{T/2}\dd t\,\ee^{\ii\omega t} = 2\sin(\omega T/2)/\omega$, tends to $2\pi\delta(\omega)$ as $T\to\infty$ by \cref{eq:C.1-sinc} with $\xi = \omega$ and $t = T/2$. The second identity is the product of three copies of the first. Inserting the inverse transform of $O(\vect{q})$ into the forward transform of \cref{eq:A.3-space} and using \cref{eq:C.1-complete} returns $O(\vect{q})$, which proves that the two transforms are inverse to each other. The transform of $\delta^3(\vect{r} - \vect{r}_0)$ is $\ee^{\ii\vect{q}\cdot\vect{r}_0}$, and the density operator \cref{eq:A.3-ne} of point particles is a sum of these phases. The transform of a derivative is $\int\dd^3r\,\ee^{\ii\vect{q}\cdot\vect{r}}\nabla O = -\ii\vect{q}\,O(\vect{q})$ after an integration by parts. Hence $-\nabla^2 \to q^2$, the step that solves the static field equations of \cref{app:C.9}.

Parseval's theorem and the convolution theorem follow from \cref{eq:C.1-complete}. Inserting the inverse transforms of $f$ and $g$ into $\int\dd^3r\,f^*(\vect{r})\,g(\vect{r})$ and doing the $\vect{r}$ integral first gives Parseval's theorem,
\begin{equation}
  \int\dd^3r\,f^*(\vect{r})\,g(\vect{r}) = \int\frac{\dd^3q}{(2\pi)^3}\,f^*(\vect{q})\,g(\vect{q}) ,
  \label{eq:C.1-parseval}
\end{equation}
and in one dimension of time $\int\dd t\,f^*(t)g(t) = \int\dd\omega\,f^*(\omega)g(\omega)/2\pi$. Let $(f\ast g)(\vect{r}) = \int\dd^3r'\,f(\vect{r} - \vect{r}')\,g(\vect{r}')$ be the convolution. Writing $\ee^{\ii\vect{q}\cdot\vect{r}} = \ee^{\ii\vect{q}\cdot(\vect{r} - \vect{r}')}\ee^{\ii\vect{q}\cdot\vect{r}'}$ and shifting the variable of the outer integral gives the convolution theorem,
\begin{equation}
  (f\ast g)(\vect{q}) = f(\vect{q})\,g(\vect{q}) .
  \label{eq:C.1-conv}
\end{equation}
The Helm form factor of \cref{sec:2.2} is an example: a uniform sphere smeared by a Gaussian is the convolution of the two densities, and the transform of the smeared sphere is the product of $3j_1(qR_{\mathrm{nuc}})/(qR_{\mathrm{nuc}})$ (\cref{app:C.3}) and the Gaussian factor $\ee^{-q^2s_{\mathrm{skin}}^2/2}$ (\cref{app:C.6}).

The time form of \cref{eq:C.1-conv} explains the structure of response functions. A causal response is a convolution in time of a response kernel with the applied field, and in frequency it becomes a product. The kernel vanishes for negative times, and its transform $\int_0^\infty\dd t\,\ee^{\ii\omega t}(\cdots)$ converges only with a factor $\ee^{-0^+t}$, a decay slower than any physical time scale. For a kernel that oscillates at a single transition energy $\Delta E = E_f - E_i$, as in \cref{app:A.3},
\begin{equation}
  -\ii\int_0^\infty\dd t\,\ee^{\ii(\omega - \Delta E)t - 0^+t} = \frac{1}{\omega - \Delta E + \ii0^+} ,
  \label{eq:C.1-causal}
\end{equation}
which is the origin of the infinitesimal $\ii0^+$ in every energy denominator of this book (\cref{app:A.3,sec:2.4}). \Cref{app:C.4} gives the meaning of the right side of \cref{eq:C.1-causal} as a distribution, and \cref{app:C.5} the analyticity that \cref{eq:C.1-causal} implies.
